% Figure — Workflow d'évaluation n8n (phase 1), vue simplifiée.
% Couleurs reprises de la figure « Génération de la télémétrie TO et TI ».
% Aucune bibliothèque TikZ supplémentaire requise (graphicx pour \resizebox).
\providecolor{couleurCorpus}{RGB}{241,247,244}
\providecolor{couleurCorpusBord}{RGB}{26,122,76}
\providecolor{couleurBleu}{RGB}{31,95,168}
\providecolor{couleurBleuFond}{RGB}{242,245,250}
\providecolor{couleurOrange}{RGB}{252,229,205}
\providecolor{couleurOrangeBord}{RGB}{214,155,80}
\providecolor{couleurN8n}{RGB}{234,75,113}
\providecolor{couleurRelance}{RGB}{184,84,80}
\resizebox{\linewidth}{!}{%
\begin{tikzpicture}[
  x=1cm, y=1cm, >=latex,
  m/.style={draw=couleurBleu, fill=couleurBleuFond, thick, rounded corners=3pt,
            align=center, text width=2.5cm, minimum height=1.9cm,
            inner sep=3pt, font=\small},
  f/.style={->, thick, draw=couleurBleu},
  retour/.style={->, thick, dashed, draw=couleurRelance},
  lab/.style={font=\footnotesize\itshape, text=couleurRelance, inner sep=2pt}
]
% ---------------- cadre n8n ----------------
\draw[draw=couleurN8n, very thick, rounded corners=8pt, fill=white] (-1.75,-2.35) rectangle (13.75,2.1);
\node[fill=couleurN8n, text=white, font=\small\bfseries, rounded corners=3pt, inner sep=4pt]
  at (0.6,2.1) {Workflow n8n};

% ---------------- modules ----------------
\node[m] (prep)  at (0,0)    {\textbf{Préparation}\\lecture du corpus\\126 appels\\(42 $\times$ 3)};
\node[m] (boucle) at (3.0,0) {\textbf{Boucle}\\un appel\\à la fois};
\node[m] (appel) at (6.0,0)  {\textbf{Appel du LLM}\\requête API\\selon le fournisseur};
\node[m] (jug)   at (9.0,0)  {\textbf{Jugement}\\extraction du JSON\\et statut};
\node[m] (ecr)   at (12.0,0) {\textbf{Écriture}\\une ligne JSON\\par exécution};

% ---------------- flux ----------------
\draw[f] (prep) -- (boucle);
\draw[f] (boucle) -- (appel);
\draw[f] (appel) -- (jug);
\draw[f] (jug) -- (ecr);
% relance si erreur transitoire
\draw[retour] (jug.north) -- (9.0,1.45) -- node[lab, above]{relance si erreur (quota, délai)} (6.0,1.45) -- (appel.north);
% appel suivant
\draw[retour] (ecr.south) -- (12.0,-1.7) -- node[lab, below]{appel suivant} (3.0,-1.7) -- (boucle.south);

% ---------------- entrée : échantillons ----------------
\node[draw=couleurCorpusBord, fill=couleurCorpus, thick, rounded corners=3pt,
      align=center, text width=2.6cm, minimum height=1.9cm, font=\small]
      (corpus) at (-4.4,0) {\textbf{42 échantillons}\\S000 à S041};
\draw[->, very thick, draw=couleurCorpusBord] (corpus.east) -- (-1.75,0);

% ---------------- sortie : fichiers JSONL ----------------
\draw[->, very thick, draw=couleurOrangeBord] (13.75,0) -- (15.45,0);
\foreach \d in {0.40, 0.20, 0}{
  \filldraw[draw=couleurOrangeBord, fill=couleurOrange, thick]
    (15.6+\d,-1.6+\d) -- (15.6+\d,1.6+\d) -- (17.6+\d,1.6+\d) -- (18.3+\d,0.9+\d) -- (18.3+\d,-1.6+\d) -- cycle;
  \filldraw[draw=couleurOrangeBord, fill=white, thick]
    (17.6+\d,1.6+\d) -- (17.6+\d,0.9+\d) -- (18.3+\d,0.9+\d);
}
\node[font=\bfseries\small, text=couleurOrangeBord] at (16.4,1.2) {JSON};
\node[font=\scriptsize\ttfamily, align=left, anchor=north west] at (15.7,0.8)
  {\{"meta": \{…\},\\\ "mapping": \{…\},\\\ "usage": \{…\}\}\\\{…\}\\\{…\}};
\node[font=\small, align=center, anchor=north] at (16.95,-1.8)
  {\textbf{Fichiers de résultats}\\\texttt{resultats\_\_<modèle>.jsonl}};
\end{tikzpicture}%
}
